\centerline{\Huge \bf  Тренировочная Олимпиада 5}
\centerline{\Huge \bf 7 класс}
\centerline{\Huge \bf Уровень: Заключительный отбор в Сириус}

\begin{flushright}
    \scriptsize{Глава 5. Вы двое, танцуйте, если хотите победить!}
\end{flushright}

\problem{1}
Найдите все группы из двух или более последовательных натуральных чисел с суммой 100
\begin{flushright}
    \textit{\small{(Гасанов Р.В.)}}
\end{flushright}

\problem{2}
На доске выписаны в ряд 5 целых чисел. Эрдни утверждает, что сумма любых двух подряд идущих отрицательных, а Соня утверждает, что суммы любых трёх подряд идущих положительна. Докажите, что кто-то из них ошибается.
\begin{flushright}
    \textit{\small{(Гасанов Р.В.)}}
\end{flushright}

\problem{3}
После очередного блока занятий в ЭлМШ несколько учеников покинули школу, а несколько наоборот присоединилось. Изначально мальчиков и девочек было поровну, но после этого блока количество учеников сократилось на 10\%, а вот доля мальчиков выросла до 55\%. Мальчиков стало больше или меньше?
\begin{flushright}
    \textit{\small{(Гасанов Р.В.)}}
\end{flushright}

\problem{4}
В треугольнике $ABC$ на стороне $AC$ выбрана точка $D$ такая, что медиана $AM$ треугольника $ABD$ параллельна медиане $DN$ треугольника $DBC$. Найдите отношение
$\dfrac{AD}{DC}$
\begin{flushright}
    \textit{\small{(Гасанов Р.В.)}}
\end{flushright}

\problem{5}
В равностороннем треугольнике $ABC$ взята точка $O$, причём $AO = a, BO = b, CO = c.$ Докажите, что существует треугольник с длинами сторон $a, b, c$ причём его вершины будут лежать на сторонах треугольника $ABC$
\begin{flushright}
    \textit{\small{(Гасанов Р.В.)}}
\end{flushright}

\problem{6}
Перед приездом РВ в Элисту Герман решил подготовиться и научиться бороться. Найдя ещё 7 единомышленников они начали тренировки. И наконец в июле решили понять, кто же из них подготовился лучше всего. Было решено провести круговой турнир, где поборются каждый с каждым по одному разу. За победу в схватке присуждаются три скибидиочка, за ничью один, за поражение 0. По итогу турнира оказалось, что те кто провели ничейную схватку оказались на разных местах. Когда АА слил РВ информацию о подготовке и турнире, РВ заметил, что число ничьих невозможно увеличить. Сколько ничьих было в турнире?
\begin{flushright}
    \vspace{-5mm}
    \textit{\small{(Очиров О.В.)}}
\end{flushright}
%https://problems.ru/view_problem_details_new.php?id=65059